\documentclass[10pt]{article}

\usepackage[utf8]{inputenc}

\usepackage[T1]{fontenc}

\usepackage{amsmath}

\usepackage{amsfonts}

\usepackage{amssymb}

\usepackage[version=4]{mhchem}

\usepackage{stmaryrd}

\usepackage{bbold}

\usepackage{CJKutf8}



\begin{document}

родном поле тяжести; понятие же о Ц. и. не связано ни с каким силовым полем и имеет смысл для любой механич. системы. Для твердого тела положения Ц. и. и центра тяжести совпадают.



При движении механич. системы ее Ц. и. движется так, как двигалась бы материальная точка, имеющая массу, равную массе системы, и находящаяся под действием всех внешних сил, приложенных к системе. Кроме того, нек-рые уравнения движения механич. системы (тела) по отношению к осям, имеющим начало в Ц. и. и движущимся вместе с Ц. и. поступательно, сохраняют тот же вид, что и для движения по отношению к инерциальной системе отсчета.



C. $M$. Tape.



ЧЕНTP КА4АНИЯ Ф и зиче ского мая тника см. Маятник.



ЦЕНТР МАСС - см. Центр инерции.



ЦЕНТР СИЛЬ - см. Центральная сила.



ЧEHTPA MACC CИCTEMA - система отсчета, в которой центр масс покоится. Конфигурационное пространство $\mathbb{R}^{3 N}$ координат $r_{i}$ системы $N$ частиц представляется в виде ортогональной суммы $\mathbb{R}_{\rho}^{3} \oplus \mathbb{R}_{x}^{3(N-1)}$, где



\[

\rho=\left(\sum_{i=1}^{N} m_{i}\right)^{-1}\left(m_{1} r_{1}+m_{2} r_{2}+\ldots+m_{N} r_{N}\right)

\]



\begin{itemize}

  \item координаты центра масс, а $x \in \mathbb{R}^{3(N-1)}$ - относительные координаты. Ц. м. с. обычно используется для описания систем нескольких частиц. Оператор энергии таких систем имеет вид:

\end{itemize}



\[

H=-\sum_{i=1}^{N}\left(2 m_{i}\right)^{-1} \Delta_{r_{i}}+V\left(r_{1}, r_{2}, \ldots, r_{N}\right)

\]



где функция $V-$ трансляционно инвариантна, то есть



$V\left(r_{1}+\lambda, r_{2}+\lambda, \ldots, r_{N}+\lambda\right)=V\left(r_{1}, r_{2}, r_{3}, \ldots, r_{N}\right), \lambda \in \mathbb{R}^{3}$.



Оператор $H$ приводится к виду



\[

H=-\left(2 \sum_{i=1}^{N} m_{i}\right)^{-1} \Delta_{\rho}+H_{\mathrm{cm}},

\]



где первое тривиальное слагаемое - оператор кинетич. энергии движения системы $N$ частиц как целого, a $H_{\mathrm{cm}}-$ оператор энергии $N$ частиц в Ц. м. с. Если в качестве координат $x$ взять набор приведенных Якоби координат, то оператор $H_{\mathrm{cm}}$ дается формулой:



\[

H_{\mathrm{cm}}=-\Delta_{x}+V(x), x \in \mathbb{R}^{3(N-1)} .

\]



Лum.: [1] Гольдбергер М., Ватсон К., Теория столкновений, пер. с англ., М., 1967; [2] Р и Д., Са йм о н Б., Методы современной математической физики, пер. с англ., т. 3, Теория $\begin{array}{ll}\text { рассеяния, М., } 1982 . & \text { С. П. Меркурьев, С.Л. Яковлев. }\end{array}$ \begin{CJK}{UTF8}{mj}山\end{CJK}EHTPATИЗАTOP подмножества $X$ в группе $G-$ множество всех $g \in G$, удовлетворяюших условию $g x=x g$ для любого $x$ из $G$. Ц. подмножества $X$ в $G$ является подгруппой в $G$. Если $A$ - ассоциативная алгебра (соответственно алгебра Ли), то Ц. подмножества $X$ в $A$ - это множество всех $a \in A$, удовлетворяющих условию $x a=a x$ (соответственно $[x, a]=0)$. Ц. подмножества $X$ в $A$ является подалгеброй в $A$. Ю. А. Неретин.



ЦЕНТРАЛЬНАЯ СИЛА - приложенная к материальному телу сила, линия действия которой при любом положении тела проходит через некоторую определенную точку, называемую ц е н т р о м с и л ы. Примеры Ц. с. - сила тяготения, направленная к центру планеты, кулоновы силы электростатич. притяжения или отталкивания и др. Под действием Ц. с. центр масс свободного тела движется по плоской кривой, а отрезок прямой, соединяющей этот центр с центром силы, описывает в любые равные промежутки времени равные плошади (см. Площадей закон).



ЧЕНТРАЛЬНОЕ МНОГООБРАЗИЕ - специальНЫМ образом выбранное инвариантное многообразие: а) векторного поля в окрестности особой точки, либо замкнутой траек- тории или б) диффеоморфизма в окрестности неподвижной точки.



Те орема 1 (см. [1]-[3], [5], [6]). Пусть $O$ - особая точка $C^{2}$-гладкого векторного поля $V$ на $N$-мерном многообразии $M$, а $V_{* 1} O^{-}$линеаризация поля $V$ в точке $O$. Пусть $T^{s}, T^{u}$ и $T^{c}$ - такие инвариантные относительно $V_{* \mid O}$ подпространства касательного пространства к $M$ в точке $O$, что собственные числа ограничения $V_{* \mid O}$ на $T^{s}, T^{u}$ и $T^{c}$ имеют отрицательные, положительные и нулевые вешественные части соответственно. Пусть размерности подпространств $T^{s}, T^{u}$ и $T^{c}$ равны $N_{s}, N_{u}$ и $N_{c}$, соответственно, $N_{s}+N_{u}+N_{c}=N$. Тогда поле $V$ имеет в окрестности точки $\boldsymbol{O}$ следующие пять проходящих через $O$ инвариантных многообразий: устойчивое (притятивающе е) $W^{s}$, не устойчивое (отталкивающе е) $W^{u}$, центральное (нейтральное) $W^{c}$, центрально-устойчивое $W^{c s}$ и центрально-неустойчивое $W^{c u}$ размерностей $N_{s}, N_{u}, N_{c}, N_{c}+N_{s}$ и $N_{c}+N_{u}$ соответственно. Многообразия $W^{s}, W^{u}$ и $W^{c}$ касаются в точке $O$ подпространств $T^{s}, T^{u}$ и $T^{c}$ соответственно. Многообразие $W^{c s}\left(W^{c u}\right)$ содержит многообразия $W^{c}$ и $W^{s}$ (соответственно $W^{c}$ и $W^{u}$ ). Решения уравнения $d x / d t=V(x)$ с начальными условиями на $W^{s}\left(W^{u}\right)$ экспоненциально стремятся к $O$ при $t \rightarrow+\infty$ (соответственно при $t \rightarrow-\infty$ ). Если поле $V$ гладкое класса $C^{r}$, где $r \geqslant 2-$ натуральное число, то многообразия $W^{s}$ и $W^{u}$ также гладкие класса $C^{r}$, а многообразия $W^{c}, W^{c s}$ и $W^{c u}$ класса $C^{r-1}$.



Если поле $V$ аналитично или $C^{\infty}$-гладко, то многообразия $W^{s}$ и $W^{u}$ также аналитичны или $C^{\infty}$-гладки соответственно, в то время как гладкость многообразий $W^{c}, W^{c s}$ и $W^{c u}$, вообще говоря, лишь конечна (но растет с уменьшением рассматриваемой окрестности точки $O$ ).



Частный случай этого утверждения, когда $N_{c}=0$, называется теоремой Адамара - Перрона (см. [5]) или теоремой Ляпунова - Перрона (см. [2]).



Если $N_{c}=0$, то $W^{c}=\{0\}, W^{c s}=W^{s}$ и $W^{c u}=W^{u}$. Многообразия $W^{s}$ и $W^{u}$ в этом случае называются се па р а т рисными многообразиями.



Теорема 1 имеет бесконечномерные аналоги (см. [4]). Кроме того, справедливы аналоги этой теоремы для неподвижных точек диффеоморфизмов (см. [1]) [в том числе (см. [4]) в бесконечномерном пространстве] и для замкнутых траекторий векторных полей (см. [1], [6]). Существуют обобщения теоремы 1 на автономные уравнения (см. [1], 12]) и на произвольные инвариантные многообразия диффеоморфизмов и векторных полей (см. [3]).



Теорема 2 (теорема сведения Шошитай ш в и л и [7]). В условиях теоремы 1 уравнение $d x / d t=V(x)$ в окрестности точки $O$ топологически эквивалентно системе



\[

d \xi / d t=-\xi, d \eta / d i=\eta, d \zeta / d t=v(\zeta),

\]



где $\xi \in T^{s}, \eta \in T^{u}$, а третье уравнение есть ограничение исходного на многообразие $W^{c}$.



Справедлив аналог теоремы 2 для неподвижных точек диффеоморфизма и ее обобщения на конечнопараметрич. семейства векторных полей и диффеоморфизмов.



Лит.: [1] Хартман Ф., Обыкновенные дифференциальные уравнения, пер. с англ., М., 1970; [2] Пл и с с В. А., Интегральные множества периодических систем дифференциальных уравнений, M., 1977; [3] Hirsch M.W., Pugh C. C., Shub M., «Lect. Notes Math.», 1977, v. 583; [4] Марсден Дж., Мак-К ракен М., Бифуркация рождения цикла и ее приложения, пер. с англ., М., 1980; [5] Арнольд В. И., Ильяшенко Ю.С., в кн.: Итоги науки и техники. Современные проблемы математики. Фундаментальные направления, т. 1, М., 1985, с. 7-149; [6] Kell a у А., «J. Diff. Equat.», 1967, v. 3, № 4, p. 546-70; [7] Шоши т а йш в ил и А. Н., «Труды семинара имени И.Г. Петровского», в. 1, М., 1975, c. $279-309$. М. Б. Севрюк.



ЧЕНТРАЛЬНЫЙ МОМЕНТ - см. Момент случайной величины.





\end{document}